\section{Related Works}

RDTU studies direct numerical text generation with supervised instruction tuning and reinforcement-driven refinement for time series forecasting. In the following, we review existing literature across two key dimensions: traditional Transformer-based approaches and emerging LLM-based methodologies for time series forecasting (TSF).


\paragraph{Transformer-based TSF}
Deep learning architectures have fundamentally reshaped time series forecasting, with Transformers initially gaining prominence due to their sequence modeling capabilities~\citep{wen2022transformers}. To address computational bottlenecks, efficient variants such as Informer~\citep{zhou2021informer}, Reformer~\citep{kitaev2020reformer}, and Pyraformer~\citep{liu2022pyraformer} introduced sparse attention mechanisms to reduce complexity for long-sequence modeling. Subsequently, decomposition-based architectures like Autoformer~\citep{wu2021autoformer} and FEDformer~\citep{zhou2022fedformer} enhanced stability on non-stationary data by explicitly disentangling seasonal and trend components. Most recently, PatchTST~\citep{nie2022time} and iTransformer~\citep{liu2023itransformer} have achieved state-of-the-art results by employing patching techniques and inverted embeddings to effectively capture multivariate correlations and local semantic contexts.


\paragraph{LLM-based TSF}

The unprecedented capabilities of foundation models have catalyzed cross-modality learning in time series, evolving from the zero-shot textual generation of LLMTime~\citep{gruver2023llmtime} to alignment-based adaptations. Notably, GPT4TS~\citep{zhou2023one} validates freezing LLM backbones with fine-tuned embeddings, while Time-LLM~\citep{jin2024timellm} employs a reprogramming framework to align temporal patches with text prototypes, and UniTime~\citep{liu2024unitime} utilizes textual instructions for cross-domain unification. Concurrently, research has expanded into visual modalities: VisionTS~\citep{chen2024visionts} adapts masked autoencoders for zero-shot forecasting by treating time series as visual signals; Time-VLM~\citep{zhong2025time} leverages vision-language models by converting data into visual plots; and DMMV-A~\citep{shen2025dmmv} enhances large vision models through multi-view constraints to optimize long-term forecasting stability. Beyond adapting general-purpose LLMs, recent studies have also explored dedicated time-series foundation models pre-trained on large-scale temporal corpora, such as Chronos-2, TimesFM, Moirai, and MOMENT~\citep{ansari2025chronos2,das2024decoder,woo2024unified,goswami2024moment}.
